\Chapter{Bevezetés}

A modern szoftverrendszerek jelentős része valamilyen adatbázis-kezelő rendszert használ az adatok tárolására és lekérdezésére. Az adatbázisok kezelésének egyik legelterjedtebb eszköze az SQL (Structured Query Language), amely deklaratív módon teszi lehetővé az adatok lekérdezését és módosítását. Az SQL széles körben alkalmazott technológia, amely mára az adatbázis-kezelő rendszerek alapvető lekérdezési nyelve.

A modern programozási környezetekben azonban gyakran szükség van arra, hogy az adat\-bá\-zis-lekérdezések szorosan illeszkedjenek az alkalmazás forráskódjához. Erre a problémára különböző programozási nyelvek eltérő megoldásokat kínálnak. A Java környezetben például a Stream API biztosít lehetőséget kollekciók deklaratív feldolgozására, míg a funkcionális programozási nyelvekben, mint például a Scala vagy a Kotlin, a beépített kollekcióműveletek hasonló szemléletet követnek. Ezek a megoldások lehetővé teszik az adatok feldolgozásának kifejezését deklaratív módon.

A .NET környezetben erre a problémára kínál megoldást a Language INtegrated Query (LINQ), amely a lekérdezési lehetőségeket közvetlenül a programozási nyelv szintaxisába integrálja. A LINQ lehetővé teszi, hogy a fejlesztők egységes módon kezeljék a különböző adatforrásokból származó adatokat, legyen szó objektumgyűjteményekről, adatbázisokról vagy XML dokumentumokról. Ennek köszönhetően az alkalmazás üzleti logikája és az adatlekérdezések ugyanabban a programozási környezetben valósíthatók meg \cite{skeet2019csharp,freeman2010prolinq}. A .NET ökoszisztémában az Entity Framework Core az egyik legelterjedtebb ORM keretrendszer, amely a LINQ lekérdezéseket közvetlenül SQL utasításokká fordítja, így a fordító célplatformjaként kézenfekvő választásnak bizonyult.

A különböző megközelítések közül jelen dokumentáció a LINQ alapú lekérdezések feldolgozásával és generálásával foglalkozik, mivel a .NET ökoszisztémában széles körben alkalmazott megoldást jelent, valamint jól vizsgálható rajta keresztül az SQL és a magasabb szintű programozási nyelvek közötti leképezés problémaköre.

A dokumentáció egy olyan fordító (transpiler) tervezését és megvalósítását mutatja be, a\-mely SQL lekérdezések automatikus LINQ kifejezésekké alakítására képes \cite{nmorph2023}. A projekt célja annak vizsgálata, hogy milyen módon valósítható meg egy ilyen fordító, milyen korlátok merülnek fel az SQL és LINQ közötti leképezés során, valamint milyen architekturális megoldások szükségesek egy ilyen rendszer létrehozásához.

Nagy nyelvi modellek (LLM) szintén képesek SQL lekérdezéseket LINQ kóddá alakítani, azonban ez a folyamat statisztikai alapú. Ugyanarra a bemenetre eltérő kimenet keletkezhet, a helyesség nem garantált, és az eredmény nem integrálható automatizált pipeline-ba. Egy determinisztikus, formálisan definiált grammatikán alapuló fordító ezzel szemben verifikálható, reprodukálható és programozottan hívható. Ezt a megközelítést vizsgálja a jelen munka.

A megvalósítás során külön figyelmet kap a lekérdezések feldolgozásának folyamata, az absztrakt szintaxisfa (Abstract Syntax Tree) alapú reprezentáció kialakítása, valamint a generált LINQ kód előállítása \cite{parr2013antlr4}.

A dolgozatban több angol eredetű szakkifejezést is használok, például a parse fa, a string, a visitor és a pipeline kifejezést. Ezeknek van magyar megfelelője (elemzési fa, karakterlánc, látogató, feldolgozási lánc), a szoftverfejlesztési gyakorlatban és az alkalmazott eszközök dokumentációjában azonban az angol alak terjedt el. A forráskóddal és a dokumentációval való egyezés érdekében ezért szándékosan ezeket használom. A formális nyelvek és a fordítóprogramok elméleti fogalmainál a magyar szakirodalom szóhasználatát követem \cite{bach2002formalis,csornyei2006fordprog}.

